% Section 9, Limitations.
\section{Limitations}
\label{sec:limitations}

Our headline result rests on a single benchmark's Python slice,
comprising 266 tasks.  While the searcher's localization ability
transfers across languages (\S\ref{sec:training}), that evidence covers
the searcher alone, not the full routed system.  Cross-benchmark
generality remains untested, and the ablation finding, that blanket
assignment to the cheapest fixer ties the routed system, is a property
of this task pool; it need not transfer to other pools or benchmarks.
The 266 tasks themselves span three repositories, with per-repository
recall ranging from 0.472 to 0.671
(Appendix~\ref{apx:localization}); the effective diversity is
closer to three codebases than to 266 independent draws.

At $n{=}266$ the headline is a one-task edge, which is precisely why we
claim a \emph{match} and nothing stronger.  The robustness evidence is
structural: both the routed system and the no-router ablation sit
$+3.60$ percentage points above the blind-mixing line, so on this pool
the margin over random traffic splitting is carried by the handoff
rather than by the routing decision.  Like the ablation finding itself,
this attribution is specific to this pool and benchmark and can change
with either.

The calibration label run uses $n{=}100$ tasks.  No handoff contrast in
this paper reaches statistical significance, on the calibration labels
or on the benchmark subset; the mechanism rests on the consistency of
the pattern across two independent settings, not on any single test.

The benchmark-side handoff comparisons themselves come from
router-selected subsets rather than random assignment, so they carry
observational caveats.  The controlled reference point is the paired
calibration study on held-out tasks (\S\ref{sec:calibration}).

Contamination control is one-sided.  The searcher's training corpus
excludes a 23-repository blocklist covering every SWE-bench Pro and
SWE-bench Verified repository, enforced by a programmatic gate on the
frozen training file.  No such control is possible on the fixer side:
the three evaluation repositories are public, we have no visibility
into the four fixers' training data, and differential fixer-side
contamination would distort the relative solve rates that the
r\'esum\'es and the calibration are fit to.

We also did not run a pass-through arm in which fixers receive the
handoff with its unverified reproduction claims intact, so the guard's
contribution to the headline solve rate is asserted from the claim
census (\S\ref{sec:results}) rather than measured; a pass-through
ablation is future work.

All costs are provider list prices recorded at measurement time.
Prices drift, and one fixer was served as a quantized build through a
pinned provider, so cost conclusions are snapshots rather than stable
constants.  The cost-accounting asymmetry described in
\S\ref{sec:setup} (all-in for the system versus API-only for solo
baselines) is conservative in our disfavor, but it is still an
accounting choice that readers should weigh.

\modelname achieves an overall localization recall of 0.566, a moderate
figure drawn from a single sampled inference pass whose variance is
uncharacterized.  The system-level result shows that this recall
suffices for effective routing; it does not establish \modelname as a
state-of-the-art localizer.

All absolute solve rates reported here are specific to the capped,
silent-budget, auto-submit evaluation tier defined in
\S\ref{sec:setup}.  Under other tiers these numbers would change.  Our
protocol comparisons (\S\ref{sec:results}) show that such tier choices
are worth percentage points of solve rate, underscoring the relativity
of any single absolute figure.

Finally, the hidden-state features that drive routing are tied to the
exact \modelname checkpoint.  Retraining the searcher would require
re-extracting hidden states and refitting the router heads.  The
zero-retrain onboarding property applies to adding new fixers, not to
changing the searcher itself.
